\documentclass[11pt]{article} 

% Paquetes de idioma y codificación
\usepackage[utf8]{inputenc}
\usepackage[T1]{fontenc}
\usepackage[spanish]{babel}

% Paquete de geometría (una sola vez)
\usepackage[
  paperwidth=21.6cm,
  paperheight=27.9cm,
  left=3cm,
  right=2cm,
  top=3cm,
  bottom=3cm
]{geometry}

% Paquetes de tipografía
\usepackage{mathptmx}
\usepackage{microtype}

% Paquetes para manejo de colores y gráficos
\usepackage{xcolor}
\usepackage{graphicx}
\usepackage{tikz}

% Definición de colores (ANTES de lstset que los usa)
\definecolor{darkblue}{rgb}{0.0, 0.0, 0.55}
\definecolor{codegreen}{rgb}{0.25, 0.49, 0.48}
\definecolor{codegray}{rgb}{0.5, 0.5, 0.5}
\definecolor{codepurple}{rgb}{0.58, 0, 0.82}
\definecolor{backcolour}{rgb}{0.95, 0.95, 0.92}

% Configuración de enlaces
\usepackage{hyperref}
\hypersetup{
    colorlinks = true,
    linkcolor  = darkblue,
    citecolor  = black,
    filecolor  = blue,
    urlcolor   = blue
}

% Paquetes para tablas y figuras
\usepackage{booktabs}
\usepackage{float}

% Personalización de códigos fuente
\usepackage{listings}
\lstset{
    literate=
    {á}{{\'a}}1 {é}{{\'e}}1 {í}{{\'i}}1 {ó}{{\'o}}1 {ú}{{\'u}}1
    {Á}{{\'A}}1 {É}{{\'E}}1 {Í}{{\'I}}1 {Ó}{{\'O}}1 {Ú}{{\'U}}1
    {ñ}{{\~n}}1 {Ñ}{{\~N}}1 {ü}{{\"u}}1 {Ü}{{\"U}}1,
    backgroundcolor=\color{backcolour},
    commentstyle=\color{codegreen},
    keywordstyle=\color{codepurple},
    numberstyle=\tiny\color{codegray},
    stringstyle=\color{red},
    basicstyle=\ttfamily\small,
    breakatwhitespace=false,
    breaklines=true,
    captionpos=b,
    keepspaces=true,
    numbers=left,
    numbersep=5pt,
    showspaces=false,
    showstringspaces=false,
    showtabs=false,
    tabsize=2,
    language=TeX,
    morecomment=[l]\#,
    frame=single,
    rulecolor=\color{black}
}

% Matemáticas y párrafos
\usepackage{amsmath}
\usepackage{parskip}
\usepackage{ragged2e}

% Configuración de secciones (una sola definición de \part)
\usepackage{titlesec}
\titleformat{\part}[display]
  {\normalfont\huge\bfseries}{}{0pt}{}
  [\thispagestyle{fancy}]

% Encabezados y pies de página
\usepackage{fancyhdr}
\pagestyle{fancy}
\fancyhf{}
\renewcommand{\headrulewidth}{0pt}
\renewcommand{\footrulewidth}{0pt}
\fancyhead[R]{\thepage}

% Bibliografía
\usepackage{natbib}
\bibliographystyle{unsrtnat}

\begin{document}

    \begin{center}

    \vspace*{1cm}

    \textbf{\Large Validez del Teorema de Cochran en la Estimación de la Varianza: 
    Un Análisis Comparativo entre Distribuciones Normal, Exponencial y t de Student}

    \vspace{1cm}

    \textbf{Lukas Wolff Casanova}\\
    Facultad de Ingeniería y Ciencias Aplicadas, Universidad de los Andes, Santiago de Chile\\
    e-mail: \href{mailto:lwolff@miuandes.cl}{lwolff@miuandes.cl}\\
    GitHub: \href{https://github.com/LukasWolff2002/TAREA_1_Estadistica_y_Dise-o}{Repositorio}
    \vspace{1cm}

    \textbf{ABSTRACT}
    \end{center}
    \vspace{0.5cm}

    This work analyzes the validity of Cochran's Theorem for variance estimation 
    across three probability distributions: Normal, Exponential, and Student's t. 
    Using samples of size $n=100$ and 1000 Monte Carlo repetitions, 95\% confidence 
    intervals for the population variance were constructed via the $\chi^2$ statistic 
    and their empirical coverage was evaluated. Results show that the Normal 
    distribution achieves exact nominal coverage ($95.5\%$), consistent with 
    Cochran's Theorem. The Exponential distribution ($95.6\%$) confirms asymptotic 
    validity through the Central Limit Theorem. However, the Student's t distribution 
    yields only $86.2\%$ coverage, revealing that heavy tails prevent convergence 
    even at $n=100$. These findings highlight that the applicability of $\chi^2$-based 
    inference depends critically on the tail behavior of the underlying distribution.

    \vspace{0.5cm}

    \textit{Key words:} \textit{Cochran's Theorem, confidence intervals, variance 
    estimation, Student's t distribution.}


\section{Introducción}

Una distribucion de probabilidad es un modelo matematico que describe como se reparte la incertidumbre asociada a una variable aleatoria. Estas distribuciones son fundamentales en estadistica, ya que permiten modelar y analizar datos, realizar inferencias y tomar decisiones informadas. En este informe, se exploraran diferentes distribuciones de probabilidad, sus propiedades y aplicaciones en el analisis de datos \cite{fisher1922mathematical}.

En términos aplicados, elegir una distribución equivale a postular una estructura para los datos: por ejemplo, la distribución Normal suele usarse cuando predominan efectos aditivos y simétricos; la Exponencial es natural para tiempos de espera y procesos de ocurrencia aleatoria; y la t de Student aparece de manera fundamental en inferencia sobre medias cuando la varianza poblacional es desconocida, especialmente en muestras pequeñas. Estas distribuciones no solo describen datos, sino que también determinan la forma de los estadísticos muestrales y, por tanto, la validez de intervalos de confianza, pruebas de hipótesis y métodos de estimación.

El presente trabajo tiene como objetivo analizar y comparar diferentes distribuciones 
de probabilidad, evaluando empíricamente la validez del Teorema de Cochran para la 
estimación de la varianza poblacional mediante intervalos de confianza basados en el 
estadístico $\chi^2$. En específico, se compararán las distribuciones Normal, 
Exponencial a tasa $\lambda = 0.5$ y t de Student con $\nu = 5$ grados de libertad. 
Además, se realizará un análisis descriptivo de la distribución t de Student, 
explicando su comportamiento y aplicaciones en el análisis de datos. Finalmente, 
se discutirán las implicancias de los resultados obtenidos y se presentarán 
conclusiones sobre la importancia de elegir adecuadamente la distribución de 
probabilidad en el análisis estadístico.

\section{Metodología}

\subsection{Teorema de Cochran}

El Teorema de Cochran \cite{cochran1934distribution} establece una propiedad fundamental de la distribución 
Normal: si $X_1, X_2, \ldots, X_n \overset{\text{iid}}{\sim} \mathcal{N}(\mu, 
\sigma^2)$, entonces la varianza muestral $S_n^2$ y la media muestral $\bar{X}$ 
son estadísticos \textit{independientes}, y además:
\begin{equation}
    \frac{(n-1)S_n^2}{\sigma^2} \sim \chi^2_{n-1}
\end{equation}

Este resultado es \textit{exacto} bajo normalidad, y constituye la base teórica 
para la construcción de intervalos de confianza para la varianza poblacional 
$\sigma^2$. En particular, un intervalo al $95\%$ de confianza se obtiene a 
partir de los cuantiles $\chi^2_{\alpha/2, n-1}$ y $\chi^2_{1-\alpha/2, n-1}$:
\begin{equation}
    \left[\frac{(n-1)S_n^2}{\chi^2_{1-\alpha/2,\, n-1}},\; 
    \frac{(n-1)S_n^2}{\chi^2_{\alpha/2,\, n-1}}\right]
\end{equation}

La condición crítica del teorema es la normalidad de los datos. Si los 
$X_i$ no provienen de una distribución Normal, la distribución exacta de 
$(n-1)S_n^2/\sigma^2$ deja de ser $\chi^2_{n-1}$, y el intervalo de confianza 
construido bajo ese supuesto puede resultar inválido. No obstante, el Teorema 
Central del Límite garantiza que, para muestras suficientemente grandes, el 
estadístico converge en distribución a $\chi^2_{n-1}$ incluso para 
distribuciones no normales, siempre que estas posean momentos finitos de orden 
suficiente y colas no demasiado pesadas.

\subsection{Distribuciones de Probabilidad}

Las distribuciones de probabilidad permiten modelar la incertidumbre asociada a variables aleatorias, describiendo cómo se distribuyen sus valores posibles. En este trabajo se consideran distintas distribuciones con el objetivo de analizar sus propiedades y su impacto en la inferencia estadística, particularmente en la estimación de la varianza poblacional.

\subsubsection{Distribución Normal}

Completamente definida por $\mu$ y $\sigma^2$, su densidad es:
\begin{equation}
f(x) = \frac{1}{\sqrt{2\pi\sigma^2}} \exp\left(-\frac{(x-\mu)^2}{2\sigma^2}\right)
\end{equation}
Es la única distribución para la cual el Teorema de Cochran aplica de forma exacta.

\subsubsection{Distribución $\chi^2$}

Surge como suma de cuadrados de variables normales estándar independientes. 
Bajo normalidad, el estadístico $(n-1)S^2/\sigma^2 \sim \chi^2_{n-1}$ permite 
construir intervalos de confianza exactos para $\sigma^2$, cuya validez fuera 
de la normalidad es lo que se evalúa en este trabajo.

\subsubsection{Distribución Exponencial}

Modela tiempos de espera con tasa $\lambda > 0$, media $1/\lambda$ y varianza 
$1/\lambda^2$. No cumple el Teorema de Cochran, pero su curtosis constante 
igual a $6$ permite que la aproximación $\chi^2$ sea asintóticamente válida 
para $n$ suficientemente grande.

\subsubsection{Distribución t de Student}

La distribución t de Student  \cite{student1908probable} surge al estimar la media de una población normal 
cuando la varianza poblacional es desconocida. Definida por sus grados de 
libertad $\nu > 0$, su densidad es:
\begin{equation}
f(x) = \frac{\Gamma\!\left(\frac{\nu+1}{2}\right)}{\sqrt{\nu\pi}\,\Gamma\!\left(\frac{\nu}{2}\right)}
\left(1+\frac{x^2}{\nu}\right)^{-\frac{\nu+1}{2}}
\end{equation}

Su varianza, definida para $\nu > 2$, se calcula como:
\begin{equation}
\operatorname{Var}(X) = \frac{\nu}{\nu - 2}
\end{equation}

lo que refleja que a mayor $\nu$, menor dispersión, convergiendo a $1$ cuando 
$\nu \to \infty$, momento en que la distribución se aproxima a la Normal estándar. 
Para $\nu > 4$, la curtosis es $\kappa = 6/(\nu-4)$, que diverge cuando 
$\nu \to 4^+$, evidenciando el peso creciente de las colas para valores pequeños 
de $\nu$.

En el mundo real, la distribución t de Student tiene aplicaciones en cualquier 
contexto donde se trabaje con muestras pequeñas y varianza desconocida: ensayos 
clínicos, control de calidad industrial y, de particular relevancia en ingeniería 
civil, en técnicas experimentales de análisis de fluidos como \textit{Particle 
Image Velocimetry} (PIV) y \textit{Particle Tracking Velocimetry} (PTV). En estos 
métodos se estiman campos de velocidad a partir de muestras reducidas en cada 
ventana de interrogación, donde la varianza del ruido de medición es desconocida. 
El estadístico $t$ permite construir intervalos de confianza sobre la velocidad 
media local y determinar si las fluctuaciones observadas son estadísticamente 
significativas o atribuibles al ruido instrumental, lo cual es crítico en flujos 
de baja turbulencia \cite{westerweel1997fundamentals, adrian2011particle}.

Respecto al Teorema de Cochran, la t de Student no satisface sus 
condiciones: al no ser Normal, el estadístico $(n-1)S_n^2/\sigma^2$ no sigue 
exactamente una distribución $\chi^2_{n-1}$. A diferencia de la Exponencial, 
cuyas colas moderadas permiten que el Teorema Central del Límite compense esta 
desviación con $n=100$, las colas pesadas de la t de Student ralentizan 
significativamente dicha convergencia, produciendo intervalos de confianza 
inválidos incluso para tamaños muestrales moderados.

\subsection{Generación de Datos Aleatorios}

Para comparar las distribuciones $\chi^2$, exponencial a tasa $\lambda = 0.5$ y t de student, se genero una muestra aleatoria de 100 datos que distrbuyen normal con $\mu = 2$ y $\sigma = 2$ con la libreria de python \texttt{scipy.stats}. Posteriormente, se realizo el siguiente procedimiento con cada una de las distribuciones mencionadas, utilizando la muestra generada:

(i) se calculo la varianza muestral ($S_n^2$) (ii) se determino un intervalo al 95\% de confianza para la varianza poblacional ($\sigma^2$) (iii) se realizaron 100 iteraciones del proceso anterior para obtener una distribucion de los intervalos de confianza y (iv) se realizaron otras 900 iteraciones para obtener un total de 1000 intervalos de confianza, con el fin de analizar la frecuencia con la que estos intervalos contienen el valor real de $\sigma^2$.

Es importante mencionar que la muestra aleatoria se generó utilizando la semilla 200, 
con el fin de garantizar la reproducibilidad de los resultados. Además, se asumen los 
siguientes supuestos: (i) los $X_i$ son independientes e idénticamente distribuidos 
(i.i.d.); (ii) la varianza poblacional $\sigma^2$ es finita, lo que impone la 
restricción $\nu > 2$ para la t de Student; (iii) el estadístico 
$(n-1)S_n^2/\sigma^2 \sim \chi^2_{n-1}$ se aplica de forma homogénea a las tres 
distribuciones, asumiendo que $n=100$ es suficiente para que el Teorema Central del 
Límite justifique la aproximación en los casos no normales; y (iv) el IC se construye 
con dos colas simétricas en probabilidad ($\alpha/2$ en cada cola), utilizando los 
cuantiles $\chi^2_{\alpha/2,\, n-1}$ y $\chi^2_{1-\alpha/2,\, n-1}$. 



\section{Resultados y Análisis}

La Figura~\ref{fig:distribuciones} muestra los histogramas de las muestras 
generadas junto a sus densidades teórica y ajustada. En el caso Normal, 
la muestra presenta un buen ajuste simétrico. La distribución Exponencial 
exhibe la asimetría positiva característica, con cola derecha extendida. 
La t de Student muestra una forma simétrica visualmente similar a la Normal, 
pero con colas más pesadas que, si bien no resultan evidentes en el histograma, 
tienen consecuencias significativas en la inferencia sobre la varianza.

\begin{figure}[h]
    \centering
    \includegraphics[width=\textwidth]{Imagenes/figura1_distribuciones.png}
    \vspace{-0.7cm}
    \caption{Histogramas de las muestras generadas con las densidades teórica 
    y ajustada para las distribuciones Normal ($\mu=2,\,\sigma=2$), 
    Exponencial ($\lambda=0.5$) y t de Student.}
    \label{fig:distribuciones}
\end{figure}


La Figura~\ref{fig:intervalos} presenta los 100 intervalos de confianza al 
$95\%$ construidos para cada distribución, ordenados de menor a mayor por 
su cota inferior. Los intervalos en azul cubren el valor real $\sigma^2$ 
(línea discontinua) y los rosados no.

\begin{figure}[h]
    \centering
    \includegraphics[width=\textwidth]{Imagenes/figura2_intervalos.png}
    \vspace{-0.7cm}
    \caption{Intervalos de confianza al $95\%$ para $\sigma^2$ en 100 
    repeticiones. Azul: cubre $\sigma^2$ real (línea discontinua); rosado: 
    no lo cubre. Cobertura en 100 repeticiones}
    \label{fig:intervalos}
\end{figure}

\textbf{Distribución Normal.} La cobertura fue de $955/1000$ ($95.5\%$), 
en plena concordancia con el nivel nominal. El Teorema de Cochran 
se cumple exactamente: bajo normalidad, $(n-1)S_n^2/\sigma^2 
\sim \chi^2_{n-1}$ de forma exacta, independientemente de los valores 
de $\mu$ y $\sigma$, por lo que el resultado no cambia al variar los 
parámetros.

\textbf{Distribución Exponencial ($\lambda = 0.5$).} La cobertura fue de 
$956/1000$ ($95.6\%$), prácticamente idéntica al caso Normal. El Teorema 
de Cochran no aplica de forma exacta, pero con $n=100$ el Teorema 
Central del Límite garantiza la convergencia suficiente de $S_n^2$, dado 
que la Exponencial tiene colas moderadas (curtosis $= 6$, invariante en 
$\lambda$). Por ello, la cobertura se mantiene estable independientemente 
del valor de $\lambda$.

\textbf{Distribución t de Student.} La cobertura fue de $862/1000$ 
($86.2\%$), muy por debajo del $95\%$ nominal: el intervalo falla casi 
el triple de las veces esperadas. El Teorema de Cochran no se 
cumple con $n=100$. La causa son las colas pesadas de la t de Student, 
cuya curtosis $\kappa = 6/(\nu-4)$ diverge para $\nu \to 4^+$, impidiendo 
que $S_n^2$ converja a $\chi^2_{99}$ con rapidez suficiente \cite{lumley2002importance}. En 
la Figura~\ref{fig:intervalos} esto se manifiesta como un sesgo sistemático: 
los intervalos fallidos se concentran en los de menor amplitud, cuya cota 
superior queda a la izquierda de $\sigma^2$. El resultado sí cambia 
con $\nu$: a mayor $\nu$ las colas se alivian y la cobertura converge al 
$95\%$; con $\nu$ cercano a $2$ el deterioro sería aún más severo.

\section{Conclusión}

Los resultados obtenidos permiten concluir que el Teorema de Cochran se cumple 
de forma exacta únicamente bajo normalidad, donde la cobertura empírica 
($95.5\%$) coincide con el nivel nominal. Para la distribución Exponencial, 
la cobertura ($95.6\%$) confirma que la aproximación $\chi^2$ es 
asintóticamente válida con $n=100$, gracias a sus colas moderadas. En 
contraste, la t de Student evidencia que colas pesadas comprometen 
severamente la validez del intervalo ($86.2\%$), incluso para muestras 
de tamaño moderado.

En cuanto a los supuestos, los resultados muestran que el más crítico es la 
normalidad de los datos: el supuesto de que $n=100$ justifica la 
aproximación asintótica resultó válido para la Exponencial, pero insuficiente 
para la t de Student. Esto implica que, antes de aplicar el estadístico 
$\chi^2$, es fundamental verificar que la distribución subyacente no presente 
colas significativamente más pesadas que la Normal.

\newpage
\bibliography{referencias}


\end{document}